\chapter{RISC-V 架构迁移设计与实现}

\section{迁移总体思路}

本课题的架构迁移工作不是简单地将原有代码逐行改写到新平台，而是在保留原系统教学逻辑的前提下，对底层体系结构相关部分进行重新适配。整体技术路线可以概括为：首先适配 Sv39 分页机制，再迁移串口、块设备和文件系统相关 I/O 路径，最后完成系统测试。

与原 IA32 版本相比，RISC-V 版本在平台依赖部分发生了较大变化。原系统依赖 x86 分页结构、PIC 中断控制器和 ATA 设备接口，而新版本面向 RISC-V 64 位平台，在分页上采用 Sv39 三级页表，在中断上采用 RISC-V trap 机制和 PLIC 外部中断控制器，在块设备上采用 VirtIO MMIO 接口。因此，迁移工作的重点不只是替换寄存器和指令，而是围绕新的硬件抽象重新建立内核底层支撑环境。

从代码编写过程看，RISC-V 迁移可以划分为以下几个阶段：第一阶段是最小可启动阶段，完成 Rust 入口、串口输出和分页前后的基本控制流切换，使内核能够在 RISC-V 平台上进入 Rust 执行环境。第二阶段是内存管理和地址空间适配阶段，完成 Sv39 页表构建、分页使能以及内核与用户空间映射。第三阶段是设备阶段，完成 UART 串口和 VirtIO 块设备驱动。第四阶段是系统功能恢复阶段，将文件系统重新接入 RISC-V 底层环境。最后是系统集成阶段，验证文件系统装载、控制台打开、init 进程创建和 shell 运行等整体功能。

\section{RISC-V 内存管理适配实现}

RISC-V 版本采用 Sv39 分页机制作为虚拟内存管理基础。Sv39 使用三级页表结构，每级页表包含 512 个表项，页大小为 4KB，虚拟地址通过多级索引转换到物理页框。与原 IA32 版本相比，Sv39 在结构上更加规则，页表项权限语义也更直接，便于在教学系统中展示虚拟地址到物理地址的层次映射关系。

当前实现中，页目录项和页表项都以 Rust 结构体表示，并通过 EntryFlags 统一封装 VALID、READ、WRITE、EXECUTE、USER、ACCESSED 和 DIRTY 等标志位。分页开启时，系统通过设置 satp 寄存器为 Sv39 模式并刷新 TLB，正式进入分页保护状态。

在地址空间布局上，RISC-V 版本同时保留了教学系统原有的高地址内核映射思想，并结合 RISC-V 平台的物理地址分布进行了重新规划。当前实现中，物理内存起始于 RAM\_BASE，内核虚拟地址从 KERNEL\_VIRT\_BASE 开始，UART、PLIC 和 VirtIO 等 MMIO 设备则统一映射到高地址 MMIO 区域，将内核代码、普通内存和设备寄存器访问清楚地区分开。

用户空间仍然位于低地址区域，并由固定数量的用户页表覆盖。代码中使用 4 个用户页表，构成统一的用户地址空间上界。MemoryDescriptor 负责管理用户进程的代码段、数据段和栈段映射，并在需要时把逻辑布局映射到实际页表中。

页表管理的初始化过程主要包括根页表建立、内核恒等映射、高地址内核映射以及用户页表数组初始化。当前实现中，系统预先维护根页目录、低地址目录、内核恒等映射目录、内核高地址目录和内核页表等静态页表结构，并在启动阶段完成清空和重新设置。内核空间映射完成后，分页通过 \texttt{satp::set} 正式启用。

对于用户空间，系统在初始化阶段预留用户页表数组，并在进程运行过程中由 MemoryDescriptor 将具体的代码段、数据段和栈段映射到对应页表项。进程切换时，系统更新当前进程用户结构所在页的映射，并刷新 TLB，以保证内核能够访问当前进程的用户态上下文。

\section{设备与 I/O 子系统迁移实现}

\subsection{UART 串口驱动适配实现}

在控制台输出和调试支持方面，RISC-V 版本不依赖早期 SBI 提供的字符输出，而是实现了基于 MMIO UART 的串口驱动。当前代码中，串口驱动直接使用 UART0\_BASE 对应的设备寄存器，通过读写接收缓冲寄存器、线路控制寄存器、FIFO 控制寄存器和线路状态寄存器完成串口初始化、发送和接收功能。

驱动初始化过程包括设置波特率、配置 8N1 格式、开启 FIFO 以及打开接收中断。发送路径通过轮询发送空闲状态写入字节，接收路径则通过检查接收就绪位从 UART 中取出输入字符。这样做的意义在于，一方面为内核启动和调试提供了稳定输出路径，另一方面也为终端输入和 shell 交互提供了底层支持。

\subsection{VirtIO 块设备驱动实现}

块设备层的迁移是本章的另一个重点。原系统在 x86 版本中依赖 ATA DMA 控制器，而 RISC-V 版本则改用 VirtIO MMIO 块设备接口。当前实现中，VirtIOBlockDriver 负责在多个 MMIO 槽位中搜索类型为块设备的 VirtIO 设备，基于 \texttt{virtio-drivers} 库完成初始化和块传输。

在数据传输过程中，驱动根据缓冲区中的块号、传输长度和读写标志调用 \texttt{read\_blocks} 或 \texttt{write\_blocks}。同时，为了适配 VirtIO 驱动框架，还实现了专门的 VirtioHal 接口，用于完成 DMA 页分配、MMIO 地址映射和共享缓冲区处理。DMA 页由内核页分配器提供，物理地址和内核虚拟地址之间的转换则通过内存管理模块统一完成。

\subsection{文件系统与设备层协同适配}

RISC-V 版本中，文件系统本身的核心结构和算法基本保持不变，但其底层块访问路径已经适配到新的设备模型之上。块设备层将请求交给 VirtIO 块设备驱动处理。对文件系统而言，底层仍然表现为统一的块设备接口。

这样设计使得文件系统、inode 读写、超级块装载和用户程序加载等逻辑可以继续复用现有结构，而设备层通过新的块设备实现接管实际 I/O 路径。RISC-V 迁移既完成了底层适配，也避免了对文件系统上层逻辑的大范围修改。

\section{本章小结}

本章围绕 Unix V6++向 RISC-V 架构迁移的设计与实现进行了分析。迁移工作以启动环境、内存管理、中断处理和设备 I/O 适配为主线，逐步完成了从 IA32 平台到 RISC-V 64 位平台的底层支撑重建。在内存管理方面，系统已适配 Sv39 分页机制并建立了新的虚拟地址空间布局；在设备层面，完成了 MMIO UART 串口驱动和 VirtIO 块设备驱动的实现；在系统集成方面，文件系统装载、控制台初始化、shell 运行均已形成完整路径。
